\section{Technical Approach}
\label{sec:approach}

This section describes the overall architecture of \sysname and the
design principles that shape it, before the next section drills into
the specific research problems that each component raises.

\subsection{Design Principles}
\label{sec:approach-principles}

Four design principles run through the \sysname architecture and
distinguish it from prior automated-assessment efforts.

\paragraph{Non-disruption as a first-class invariant.}
Every agent that touches the real \cps is strictly passive. Active
interaction---probing, fuzzing, exploit validation, even latency-sensitive
queries---is performed exclusively against the digital twin. This
invariant is enforced at the architectural level, not left to operator
convention: the Observer Agents are given capabilities that are
passive by construction (packet capture, log reading, configuration
polling), and all action-capable agent roles run in an environment
whose only reachable network is the twin's.

\paragraph{Agents as cooperating specialists, not monoliths.}
Rather than building one large agent that attempts to do everything,
\sysname decomposes the workflow across four roles, each with a tight
specification, a narrow tool set, and a small surface of interaction
with the others. This improves reliability (failure modes are
localized), cost (specialist agents can use smaller models where
appropriate), and security (the principle of least privilege can be
applied to agents, not just users).

\paragraph{Knowledge as the shared spine.}
The KB is not a log of what the system did; it is a living model of
what the system \emph{is}. Every agent reads from it, writes to it, and
treats it as the authoritative source of truth. Contradictions between
documentation and observation are recorded as first-class facts, not
resolved silently, so that the human operator can be brought in when
the agents cannot reconcile them.

\paragraph{Closed-loop operation.}
\sysname is not a pipeline that runs end-to-end once. It is a set of
feedback loops: documentation updates the KB, which updates the
observation plan, which updates the KB, which updates the twin, which
exposes weaknesses, which produce analyst decisions, which update the
KB. Every component is designed to be incremental: the KB supports
additions and revisions, the observation plan supports live
reconfiguration, the twin supports partial reconstruction rather than
only full rebuilds, and the red team supports re-running only the
analyses whose assumptions have changed.

\subsection{System Architecture}
\label{sec:approach-arch}

Figure~\ref{fig:arch} shows how the four agent teams, the real \cps,
the digital twin, and the human analyst are connected.

\begin{figure}[t]
\centering
\begin{tikzpicture}[
    font=\small,
    node distance=6mm and 9mm,
    agent/.style={rectangle, rounded corners=2pt, draw=black!70, fill=blue!6,
                  minimum width=26mm, minimum height=8mm, align=center,
                  inner sep=3pt},
    store/.style={cylinder, shape border rotate=90, aspect=0.25, draw=black!70,
                  fill=gray!10, minimum width=22mm, minimum height=10mm,
                  align=center, inner sep=2pt},
    system/.style={rectangle, draw=black!70, fill=green!6, minimum width=30mm,
                   minimum height=10mm, align=center, inner sep=3pt},
    twin/.style={rectangle, draw=black!70, dashed, fill=orange!6,
                 minimum width=30mm, minimum height=10mm, align=center,
                 inner sep=3pt},
    human/.style={rectangle, draw=black!70, fill=red!6, minimum width=26mm,
                  minimum height=8mm, align=center, inner sep=3pt},
    arr/.style={-{Latex[length=2mm]}, thick, draw=black!65},
    darr/.style={{Latex[length=2mm]}-{Latex[length=2mm]}, thick, draw=black!65},
  ]

  % KB in the center
  \node[store] (kb) {Knowledge\\Base (Graph + Vector)};

  % Agent teams arranged around the KB
  \node[agent, above left=10mm and 2mm of kb]  (kbA)  {KB Agents};
  \node[agent, above right=10mm and 2mm of kb] (obsA) {Observer Agents};
  \node[agent, below left=10mm and 2mm of kb]  (repA) {Replication Agents};
  \node[agent, below right=10mm and 2mm of kb] (rtA)  {Red Team Agents};

  % External elements
  \node[store, left=14mm of kbA] (docs) {Documentation\\Policies, IaC};
  \node[system, below=10mm of docs] (cps) {Real Cyberphysical\\System};
  \node[twin, right=14mm of obsA] (twin) {Digital Twin\\(VMs, containers,\\emulators)};
  \node[human, below=18mm of kb] (human) {Human Analyst};

  % Arrows - data flow into KB
  \draw[arr] (docs) -- (kbA);
  \draw[arr] (cps.east) -- node[above,font=\scriptsize,pos=0.45]{passive obs.} (obsA.west);
  \draw[arr] (kbA) -- (kb);
  \draw[arr] (obsA) -- (kb);

  % Arrows - KB drives replication and red team
  \draw[arr] (kb) -- (repA);
  \draw[arr] (kb) -- (rtA);
  \draw[arr] (repA.east) -- node[above,font=\scriptsize,pos=0.3]{build} (twin.south west);
  \draw[arr] (rtA.east) -- node[above,font=\scriptsize,pos=0.3]{attack} (twin.south east);

  % Arrows - human in the loop
  \draw[arr] (rtA) -- node[right,font=\scriptsize]{findings} (human);
  \draw[arr] (human.west) to[bend right=18] node[below,font=\scriptsize,pos=0.55]{validation} (cps.south east);
  \draw[arr] (human.north) -- node[right,font=\scriptsize,pos=0.5]{feedback} (kb.south);

  % Plan coordination between KB and Observer agents
  \draw[darr, dashed] (kbA) -- node[above,font=\scriptsize]{plan update} (obsA);

\end{tikzpicture}
\caption{\sysname architecture. Four agent teams (KB, Observer,
Replication, Red Team) cooperate around a hybrid knowledge base. Only
the Observer Agents touch the real \cps, and they do so passively. All
attack activity occurs against the digital twin. Human analysts
validate findings on the real system and feed both confirmations and
false-positive diagnoses back into the knowledge base.}
\label{fig:arch}
\end{figure}

\subsection{Agent Roles in Detail}
\label{sec:approach-roles}

\subsubsection{Knowledge Base Agents}
KB Agents are responsible for the long-lived representation of the
target \cps. Their inputs are heterogeneous: vendor manuals in PDF,
network diagrams in Visio or Draw.io, configuration files in YAML and
INI, infrastructure-as-code in Terraform or Ansible, software bills of
materials, engineering change logs, and, critically, the operator's
security policy. KB Agents use LLMs to extract structured facts from
unstructured inputs, write those facts into a graph database, and
retain the source text (with provenance) in a vector store for later
retrieval. When inputs conflict, or when a required fact is missing
(e.g., the network diagram references a PLC whose firmware version is
not documented), the KB Agents generate a targeted question for a
human operator rather than guessing.

\subsubsection{Observer Agents}
Observer Agents translate the KB's picture of the system into a
concrete observation plan. The plan specifies, for each subnet of
interest, what passive vantage point will be used (SPAN port, TAP,
service-mesh sidecar, host-level eBPF probe, log shipper), what
signals will be collected, what volume of traffic will be retained,
and what privacy or safety constraints apply. Observer Agents deploy
themselves (typically as lightweight containers in an adjacent
segment) and continuously feed the KB with three kinds of update:
\emph{confirmations} of previously-known facts (which raise confidence
scores), \emph{discoveries} of previously-unknown facts (which create
new nodes in the graph), and \emph{contradictions} with documentation
(which are flagged for human attention).

\subsubsection{Replication Agents}
Replication Agents are the bridge between what the system \emph{is}
(as described by the KB) and what the defender can \emph{attack} (a
digital twin). For each class of component in the KB, a Replication
Agent selects a replication strategy: a full virtual machine for
commodity servers, a container for microservices, an emulator such as
QEMU for a non-x86 controller, an abstraction-layer rehosted firmware
image (as in HALucinator) for an embedded device whose hardware we
cannot fully emulate, or a behavioral \emph{mock} (as in Conware) for
a component we can observe but cannot run. The agent then assembles
these replicated components into a network topology that mirrors the
real one, populates it with representative data, and exposes it to
the Red Team Agents. The controlled-fidelity aspect of this process is
discussed in detail as Thrust~2 (\S\ref{sec:thrust2}).

\subsubsection{Red Team Agents}
Red Team Agents operate exclusively against the twin. They instantiate
a clean snapshot, select a ``flag'' whose compromise represents a
security-meaningful outcome (e.g., unauthorized modification of a
setpoint, escalation to an engineering workstation, exfiltration of
historian data), and plan a multi-step attack from an initial attacker
position. Atomic steps---individual exploits, misconfigurations,
privilege escalations---are drawn from a library that is continually
updated with findings from component-level analyses (fuzzing, static
analysis, symbolic execution) applied to the twin's software. Steps
are composed into a plan via PDDL, using the precondition/postcondition
formulation we developed in ChainReactor. Successful plans are
replayed on a clean twin to eliminate side-effect confounds, and are
then promoted to findings.

\subsection{Human-in-the-Loop Design}
\label{sec:approach-human}

\sysname is not a fully autonomous system and does not aspire to be
one. The human analyst is the arbiter of ground truth: they decide
whether a twin-reproducible attack is realistic, whether to attempt
validation on the real system, and whether a false positive reflects a
policy misalignment or a twin-fidelity bug. Three design choices
support this role. First, every finding carries a \emph{provenance
trace} through the KB, the observations, and the twin, so that the
analyst can see why the system believes the attack is real. Second,
the analyst's decisions are first-class inputs to the system: a
``false positive due to fidelity gap'' verdict triggers a specific
remediation workflow in the KB and Replication Agents. Third, the
validation process on the real system is itself scaffolded by the
agents, which propose a minimally invasive validation procedure that
the analyst can approve, modify, or reject.
